\documentclass[aps,prd,reprint,superscriptaddress,nofootinbib,longbibliography,floatfix]{revtex4-2}
\usepackage[T1]{fontenc}
\usepackage{amsmath,amssymb,bm}
\usepackage{graphicx,booktabs}
\usepackage[colorlinks=true,linkcolor=blue,citecolor=blue,urlcolor=blue]{hyperref}
\graphicspath{{figures/}}
\input{generated/numbers}
\newcommand{\mfl}{m_{4\ell}}
\newcommand{\W}{W_{68}}
\newcommand{\GeV}{\,\mathrm{GeV}}
\newcommand{\Pois}{\operatorname{Pois}}
\newcommand{\med}{\operatorname{median}}
\newcommand{\Mempty}{M_{\varnothing}}
\begin{document}

\title{Kinematic feature attribution for signal-strength inference in simulated
\texorpdfstring{$H\to ZZ^*\to2e2\mu$}{H to ZZ* to 2e2mu} events}
\author{[Author name to be supplied]}
\email{[Corresponding-author email to be supplied]}
\affiliation{[Institution and postal address to be supplied]}
\date{September 22, 2026}

\begin{abstract}
The usefulness of a kinematic representation for a Higgs analysis depends on
the precision of the subsequent parameter inference, as well as on its ability
to separate signal from background. We investigate this distinction in a
controlled Monte Carlo study of $H\to ZZ^*\to2e2\mu$. Nineteen observables are
partitioned into four physically motivated groups, and all fifteen nonempty
combinations are evaluated using five independently initialized classifiers
per combination. Explicit four-lepton mass is excluded from the classifiers.
A constant-score reference and a common, frozen likelihood define a complete
coalition comparison based on the expected 68\% profile-likelihood interval
width for the signal strength. The executed analysis has one mass bin spanning
105--140 GeV and two score categories for the nonempty combinations. At an
injected signal strength of unity and a nominal normalization of
$10\,\mathrm{fb}^{-1}$, the six-variable BC and ten-variable AC combinations
have median widths \WidthBC{} and \WidthAC{}, compared with \WidthABCD{} for all
nineteen variables. Their median reductions relative to the constant reference
are \GainBC\% and \GainAC\%. Exact Shapley decomposition assigns the largest
positive contributions to dilepton and collective kinematics, whereas the
angular group has a negative contribution within this procedure. Classification
and inference rankings differ. These findings are exploratory: only 161 of 200
event-group bootstrap replicas are valid, the assessment review is not
independent, and pseudo-experiment comparisons use an artificial coupling.
The results identify compact candidates for independent study, without
establishing general precision gains or validated interval coverage.
\end{abstract}
\maketitle

\section{Introduction}
\label{sec:introduction}

The decay of a Higgs boson into four charged leptons provides a particularly
transparent setting in which to study the relation between event kinematics
and statistical inference. The final state is fully reconstructible to a good
approximation, and its invariant mass is a direct handle on the resonance.
The two dilepton systems also encode the geometry of the decay, while their
motion in the laboratory reflects production dynamics and acceptance.
Experimental studies combine these ingredients to constrain the Higgs signal
strength and other properties\cite{CMS2017}. Matrix-element discriminants
illustrate the usefulness of the multidimensional four-lepton configuration
for distinguishing Higgs decays from continuum production\cite{Avery2013}.
An analysis of learned representations in this channel can therefore connect
a concrete physical description to a well-defined inference target.

That connection is not automatic. A classifier is usually trained to separate
two event populations, whereas an analysis ultimately estimates a parameter
through a likelihood. A better ranking of signal and background events need
not produce a narrower confidence interval after score compression, category
construction, and nuisance-parameter profiling. The area under the receiver
operating characteristic curve (AUC) integrates discrimination across many
thresholds; the final fit uses a particular partition with particular event
yields and statistical uncertainties. Methods such as INFERNO explicitly
address this distinction by optimizing an inference objective during
training\cite{INFERNO2019}. Here the learner remains a conventional classifier.
The question is instead how different physically motivated inputs perform
when passed through the same prescribed learning and inference procedure.

A second distinction concerns information and its representation. Derived
kinematic variables can make a useful pattern easier for a finite network to
learn without adding independent event information. Conversely, a group with
clear physical content can degrade the performance of a particular fitted
model when simulation is limited. Training fluctuations, overlapping
observables, the choice of loss, and coarse score categories all intervene
between the underlying probability distributions and the resulting interval.
Consequently, a numerical attribution to an input group should be interpreted
as an attribution of performance for a specified procedure. It is not a
measurement of intrinsic physical information.

We study this problem by partitioning nineteen reconstructed observables into
four groups and evaluating all nonempty group combinations. The four-lepton
mass is omitted as an explicit classifier input, allowing the study to focus
on the remaining representation while keeping the selected mass window and
the likelihood prescription common. For each of five training seeds, the
fifteen trained models are complemented by a constant-score reference. This
complete set permits an exact Shapley decomposition of the change in expected
signal-strength interval width\cite{Shapley1953}. Direct subset comparisons
are retained alongside the attribution, because an ordering of Shapley values
does not by itself identify an optimal feature subset\cite{Fryer2021}.

The present analysis uses only two simulated members of the ATLAS 2020 open
data collection\cite{ATLAS2020}. It is a methodological study of a restricted
population, rather than an experimental measurement or an ATLAS or CMS
collaboration result. The analysis yields a complete nominal comparison and
conditional pseudo-experiment diagnostics, but its finite-simulation and
independent-validation requirements are not all satisfied. We therefore
report the observed pattern and its limitations together. In particular,
the common mass grid has merged to a single bin. The resulting comparison
quantifies the value of score categories in that coarse analysis; it cannot
establish an improvement over a fit that resolves the Higgs mass peak.

\section{Simulated samples and event definition}
\label{sec:samples}

\subsection{Population and reconstruction}

The input collection contains proton--proton collision simulations at
$\sqrt{s}=13$ TeV. We use the signal member with dataset identifier 345060,
labelled \texttt{ggH125\_ZZ4lep}, and the background member with identifier
363490, labelled \texttt{llll}. The former defines the signal scaled by the
parameter $\mu$ throughout this paper. The latter defines the complete
background of this restricted model. Its sample label alone does not justify
identifying it with a pure $q\bar q\to ZZ$ component. Other Higgs production
modes and additional experimental backgrounds are not included. The dataset
release and individual member identities are fixed, and no collision-data
events are used.

The selected final state contains exactly two electrons and two muons, with
four-lepton mass $105\leq\mfl<140\GeV$. Events must satisfy the electron or
muon trigger, contain exactly four good leptons after quality requirements,
and have at least one good lepton matched to a trigger object. Good leptons
pass the stored tight-identification decision. Both track and calorimeter
isolation divided by lepton transverse momentum must be below 0.3. The
transverse impact-parameter significance satisfies $|d_0/\sigma_{d_0}|<5$
for electrons and $<3$ for muons; the longitudinal requirement is
$|z_0/\cosh\eta|<0.5$ mm. Electron and muon acceptance is restricted to
$|\eta|<2.47$ and $|\eta|<2.7$, respectively.

Leptons are ordered by decreasing transverse momentum, with thresholds of
20, 15, 10, and 7 GeV. The total charge must vanish, and two same-flavour,
opposite-sign pairs must be available. The pair closest to the nominal
$Z$ mass is called $Z_1$; the remaining pair is called $Z_2$. Every
same-flavour, opposite-sign pair must have mass above 5 GeV, with the chosen
pairs additionally satisfying $50<m_{Z_1}<106\GeV$ and
$12<m_{Z_2}<115\GeV$. The mixed-flavour restriction reduces the ambiguity
associated with pairing identical leptons. It also restricts the reach of the
study: conclusions for $4e$ and $4\mu$, where exchange effects and pairing
choices require further care, do not follow automatically.

These selections define a fixed analysis population. Their thresholds are
not optimized using the reported subset ranking. Reconstructed four-vectors
and the stored lepton ordering determine the observables used below; masses
and momenta are expressed in GeV, and azimuthal angles in radians. Invalid
four-vectors, undefined angular configurations, and nonfinite feature values
are treated as analysis failures rather than assigned artificial replacement
values. A complete generator, shower, and detector-provenance audit remains
necessary before using the population to support broader physical claims.

\subsection{Weights and statistical roles}

The nominal luminosity is $L=10\,\mathrm{fb}^{-1}$. For a generated event $i$,
the normalization has the form
\begin{equation}
 w_i=L\,\sigma\,k\,\epsilon_{\mathrm f}
          \frac{w_i^{\mathrm{gen}}}{\sum_j w_j^{\mathrm{gen}}},
 \label{eq:weight}
\end{equation}
where the cross section, correction factor, and filter efficiency are those
bound to the sample configuration. The physical sign of the generator weight
is retained. The study does not infer the origin of negative weights from
their presence alone: subtraction, matching, and interference are distinct
possibilities. Their implications for the template-statistical model must be
assessed separately\cite{Mandrik2017,Gluesenkamp2018}.

An existing development/test partition is preserved. The historical
development population has selection probability $p_{\mathrm{dev}}=0.8$.
Within it, physical event groups are assigned deterministically to training,
validation, calibration, template construction, and assessment, with fractions
0.4, 0.1, 0.2, 0.2, and 0.1. A group is identified by the simulation channel
and event number, so repeated representations of the same physical event
cannot cross roles. For a role with conditional probability $r$, its yield
weight is $\widetilde w_i=w_i/(p_{\mathrm{dev}}r)$. This sampling correction
is applied to yields and squared consistently in the variance; role totals
are not forced to agree by subsequent renormalization.

The roles separate decisions with different statistical functions.
Training determines network parameters and the feature standardization.
Validation selects a checkpoint. Calibration fixes score thresholds, while
the template role estimates the yields used by the likelihood. Assessment
provides a separate parent for post-freeze diagnostics. This separation
limits immediate feedback between fitting and evaluation, but it cannot
erase historical access to a sample. The retained analysis has only a
single-researcher, non-independent access review. Its assessment diagnostics
are consequently not treated as independent confirmation.

Table~\ref{tab:setup} summarizes the executed configuration. The template
population contains 7,221 signal groups and 590 background groups. Under the
bound normalization, their total expected yields are 2.76832 and 2.49785,
respectively. These small expectations are relevant to the broad
signal-strength intervals and to departures from simple large-sample
intuition. They are predictions of the restricted normalization and selection
used here, not an estimate of the total experimental yield in this channel.

\begin{table}[tbp]
\caption{Samples and fixed analysis configuration. Group counts refer to the
published non-assessment audit. Expected yields include the role-sampling
correction. The initial fine mass grid and the executed grid are distinguished.}
\label{tab:setup}
\begin{ruledtabular}
\begin{tabular}{ll}
Quantity & Configuration \\
\hline
Signal / background DSID & 345060 / 363490 \\
Collision energy & 13 TeV \\
Final state; luminosity & $2e2\mu$; $10\,\mathrm{fb}^{-1}$ \\
Mass window & $105\leq\mfl<140\GeV$ \\
Training groups ($s/b$) & 14,377 / 1,140 \\
Validation groups ($s/b$) & 3,503 / 312 \\
Calibration groups ($s/b$) & 7,262 / 575 \\
Template groups ($s/b$) & 7,221 / 590 \\
Template yields ($s/b$) & 2.76832 / 2.49785 \\
Initial / executed mass bins & 35 / 1 \\
Score categories, nonempty subsets & 2 \\
Hidden-layer widths & 64, 64, 32 \\
Training seeds & 42--46 \\
Nominal inference & T1, Asimov, $\mu=1$ \\
\end{tabular}
\end{ruledtabular}
\end{table}

\section{Kinematic representations and learning procedure}
\label{sec:representations}

\subsection{Feature groups and the mass convention}

Table~\ref{tab:features} defines the four groups. Group A contains individual
lepton transverse momenta and pseudorapidities, ordered by transverse
momentum. Group B describes the masses and internal angular separation of
the dilepton pairs. Group C contains collective four-lepton transverse
momentum and the azimuthal separation of the two dilepton systems. Group D
contains five angular coordinates associated with the production and decay
geometry. The scalar products and rest-frame constructions underlying these
angles are useful descriptions of the four-lepton phase space, as illustrated
by matrix-element studies\cite{Avery2013}. The present group definition is a
choice of representation for a classifier, not an assertion that the four
groups form independent components of phase space.

The angular convention fixes the charge ordering and beam direction.
In the four-lepton rest frame, $\theta^*$ is the angle between $Z_1$ and
the boosted positive laboratory beam direction. The decay-plane normals
are proportional to $\bm p_{\ell^-}\times\bm p_{\ell^+}$ for each pair
in that frame. The production-plane normal is proportional to the cross
product of the boosted beam direction with the $Z_1$ direction.
The two azimuthal coordinates are oriented angles about $Z_1$: first from
the $Z_1$ decay-plane normal to the $Z_2$ normal, and second from the
production-plane normal to the $Z_1$ normal. Both lie in $[-\pi,\pi)$.
For $\theta_1$ ($\theta_2$), the negative lepton in the $Z_1$ ($Z_2$)
rest frame is compared with the direction opposite the other dilepton
system. These conventions specify the implemented coordinates, without
implying an independently validated matrix-element backend.

\begin{table}[tbp]
\caption{Classifier inputs. Here $i=1,\ldots,4$ denotes leptons ordered by
decreasing $p_T$, and $\Delta R=\sqrt{(\Delta\eta)^2+(\Delta\phi)^2}$ is
evaluated within a dilepton pair. Explicit $\mfl$ is absent from every group.}
\label{tab:features}
\begin{ruledtabular}
\begin{tabular}{clc}
Group & Observables & Count \\
\hline
A & $p_{T,i},\eta_i$ & 8 \\
B & $m_{Z_1},m_{Z_2},\Delta R_{Z_1},\Delta R_{Z_2}$ & 4 \\
C & $p_{T,4\ell},\Delta\phi_{ZZ}$ & 2 \\
D & $\cos\theta^*,\cos\theta_1,\cos\theta_2,$ & \\
  & $\phi_{\mathrm{decay}},\phi_{\mathrm{production}}$ & 5 \\
\end{tabular}
\end{ruledtabular}
\end{table}

Removing the explicit mass column does not imply statistical independence
from mass. For approximately massless leptons, for example,
\begin{equation}
 m_{ij}^{2}=2p_{T,i}p_{T,j}
    \bigl[\cosh(\eta_i-\eta_j)-\cos(\phi_i-\phi_j)\bigr].
 \label{eq:pairmass}
\end{equation}
Correlated kinematic quantities can therefore retain information related to
invariant masses. Some groups omit ingredients needed for an exact
reconstruction, yet can still carry substantial mass dependence. Likewise,
the selected mass window already conditions the population before training.
We use the phrase ``without explicit mass input'' literally, and do not call
the learned scores mass-decorrelated or mass-independent.

The complete combination set consists of the fifteen nonempty subsets of
$N=\{A,B,C,D\}$. Each subset is trained separately. The BC and AC inputs
therefore contain six and ten variables, respectively, while ABCD contains
nineteen. The attribution analysis reuses the resulting 75 fixed networks;
it does not remove columns from a network trained on all variables. Removing
columns from an existing network would define a different intervention,
confounding the representation comparison with the response of a model to
inputs outside its training configuration.

\subsection{Network fitting and the empty coalition}

The learner is a multilayer perceptron with hidden widths 64, 64, and 32.
Each hidden affine transformation is followed by layer normalization and a
SiLU activation. Dropout with probability 0.1 follows the first two hidden
activations, and a final affine layer produces a single logit. Inputs are
standardized using means and population standard deviations fitted only on
the training role. A zero scale is replaced by one. The resulting
transformation is kept fixed for validation, calibration, and inference.

Training uses binary cross entropy, CPU single-precision arithmetic, and
AdamW with learning rate $10^{-3}$, weight decay $10^{-4}$, and batch size
1,024. The event weight entering optimization is the absolute physical weight
divided by its training-class mean. This provides a nonnegative training
measure while preserving the signed weights for physical yields. In
particular, the fitted logit should not be identified without qualification
with the likelihood ratio of the signed physical template distributions.
The optimizer sees a surrogate measure; the downstream likelihood determines
the reported inference performance.

Ordinary training lasts at most 200 epochs. The selected checkpoint maximizes
validation AUC under absolute physical weights, using patience 20 and minimum
improvement $10^{-4}$. The five seed values change initialization and
training randomness, while the event population and statistical roles remain
fixed. Equal seed labels provide a controlled comparison convention across
subsets, not identical optimization trajectories: changing input dimension
also changes the first layer and its parameter count. The calculation tests
this fixed learner and training prescription, rather than a capacity-matched
family or a search over architectures.

The empty coalition is a deterministic constant score, denoted
$\Mempty$ and represented by the run identity \texttt{M0off}. It receives no
kinematic inputs and has no fitted network. All events share its score and
therefore occupy one score category; the other category is structurally
empty. Ties are not split by class labels or event identifiers. The reference
is evaluated with the same accepted population, mass boundaries, and
likelihood prescription as the nonempty combinations. Five identities are
retained to align the comparison bookkeeping, but their numerical widths
are identical and do not represent five independent baseline fits.

This empty coalition makes the meaning of attribution explicit. A positive
contribution is a reduction in interval width relative to combinations with
less input information under the specified retraining procedure. It is not
an explanation of an individual event score. The coalition value includes
learning, threshold selection, compression, and template statistics, so its
sign need not match an interpretation based only on the discriminating power
of an ideal continuous observable.

\section{Statistical inference and attribution}
\label{sec:inference}

\subsection{Score categories and the common template}

The sigmoid of the network logit defines its raw score. Scores are divided
into two categories using a background equal-yield threshold estimated from
the calibration role. Every accepted
event belongs to a category: no preliminary score cut removes low-scoring
events. The calibration and template populations are distinct, so the
background yields in the two template categories need not be exactly equal.
The raw-score analysis used here does not apply a conditional cumulative
distribution transform or adversarial mass decorrelation.

The signed calibration weights require a defined prescription for that
threshold. Scores are first accumulated in twenty equal bins on $[0,1]$.
Their signed yields are fitted to a nonnegative histogram with the same
total yield by minimizing the variance-weighted squared adjustment. Bins
with zero variance are treated as structural zeros. The median of the
resulting cumulative histogram, interpolated linearly within its containing
score bin, sets the category boundary. Equal scores are assigned together,
with equality at the threshold assigned to the higher category. This
constrained distribution fit determines the threshold; it is not a
conditional transformation of the event score. In particular, the physical
template yields used by the likelihood are not replaced by the fitted
calibration histogram.

The analysis initially specifies 1 GeV mass bins over the selected window.
Support requirements are imposed on signed template yields, effective
counts, and cancellation. The effective-count threshold is 20 and the
minimum ratio $|\sum_i w_i|/\sum_i|w_i|$ is 0.2. A common deterministic merging rule is applied
before freezing, rather than optimizing a separate grid for each candidate.
For the executed analysis this procedure leaves the two edges
$[105,140]\GeV$, giving one mass bin shared by all candidates and training
seeds. Each nonempty classifier is thus evaluated through two score counts,
while the constant reference reduces to a single accepted-event count.

This coarse configuration is central to the interpretation. The likelihood
still has a registered mass coordinate, but within the accepted window it
does not resolve the peak shape. The observed gain over the constant
reference consequently reflects the extra separation supplied by the score
partition, including any residual mass dependence that its inputs encode.
It is not the gain of kinematics after exploiting a finely resolved mass
spectrum. A comparison with such a spectrum would require a distinct,
adequately supported analysis with its own fixed binning and validation.

The seed-42 BC template gives a concrete example of the compression.
Its two background category expectations are approximately 1.26935 and
1.22850 events, while the signal expectations are 0.23310 and 2.53522.
The category totals reproduce the common yields in Table~\ref{tab:setup}.
The score partition therefore concentrates most of the simulated signal
in one category while leaving roughly half the background in each.
This example explains how categorization can help the count-based inference
even when no event is rejected. It does not establish the stability of
the same threshold or category yields under new simulated samples.

\subsection{Likelihood and interval construction}

Let $s_{cj}$ and $b_{cj}$ denote signal and background expectations in score
category $c$ and mass bin $j$. The fitted model is
\begin{align}
 \mathcal L(\mu,\bm\gamma)={}&
 \prod_{c,j}\Pois\!\left(n_{cj}\mid\lambda_{cj}\right)
 \prod_{p,c,j}\Pois\!\left(a_{p,cj}\mid
       \tau_{p,cj}\gamma_{p,cj}\right),\nonumber\\
 \lambda_{cj}={}&\mu s_{cj}\gamma_{s,cj}+b_{cj}\gamma_{b,cj},
 \label{eq:likelihood}
\end{align}
where $p$ labels the signal and background processes. The auxiliary
constraints describe a finite-template approximation, called T1. For a
positive nominal yield $y$ and estimated variance $V$,
\begin{equation}
 \tau=\frac{y^2}{V},\qquad
 y=\sum_g W_g,\qquad V=\sum_g W_g^2,
 \label{eq:tau}
\end{equation}
with $W_g$ the total yield contribution of physical group $g$ to that bin.
The multiplicative nuisance parameter $\gamma$ is nonnegative, and the
nominal auxiliary expectation is $\tau$. The implementation uses the
\texttt{shapesys} modifier of pyhf 0.7.6\cite{pyhf2021,pyhf076}, with
independent process/bin auxiliary constraints.

Finite simulation must be distinguished from fixed-template counting
uncertainty\cite{Barlow1993}. Equation~\eqref{eq:tau} is an effective-count
approximation for weighted templates, not an exact generative construction
for arbitrary signed samples. Cancellation and physical-group correlations
can invalidate a naive independent-bin treatment. Nonpositive rates are not
repaired with absolute values or small positive offsets, and unsupported
configurations are recorded as failures. The fixed-template limit, termed
T0, omits the finite-template nuisance parameters. The nominal comparisons
in this paper use T1 throughout; neither its implementation nor its numerical
convergence establishes independent validity for the present signed sample.

The simpler T0 model helps explain why classification and inference can
rank representations differently. For fixed, known category expectations
and one mass bin, the Fisher information for $\mu$ is
\begin{equation}
 \mathcal I_{\mathrm{T0}}(\mu)=\sum_c
                 \frac{s_c^2}{\mu s_c+b_c}.
 \label{eq:fisher}
\end{equation}
Categories are valuable when they separate regions with different
signal-to-background composition. By the Cauchy--Schwarz inequality,
the expression is at least $s_{\mathrm{tot}}^2/
(\mu s_{\mathrm{tot}}+b_{\mathrm{tot}})$, the information in the summed
count, with equality when their compositions provide no useful
distinction. This is a local, ideal fixed-template statement. It provides
physical intuition for score categorization, not an alternative numerical
result for the finite-sample profile intervals reported here.

In particular, changing feature inputs does not simply refine a fixed
category partition: retraining moves its boundary and can produce a less
favourable partition. T1 also profiles nuisance parameters whose constraints
depend on the simulation populating each category. The ideal information
inequality therefore does not require the learned ABCD partition to
outperform BC, nor does it guarantee that the expected finite-sample
interval follows the inverse square root of Eq.~\eqref{eq:fisher}.
The distinction separates the statistical value of retaining exact event
information from the performance of a learned and discretized summary.

The profile-likelihood statistic is
\begin{equation}
 q(\mu)=-2\log
 \frac{\mathcal L(\mu,\widehat{\widehat{\bm\gamma}}(\mu))}
      {\mathcal L(\widehat\mu,\widehat{\bm\gamma})}.
 \label{eq:profile}
\end{equation}
We invert it using the one-degree-of-freedom chi-squared quantile for the
chosen confidence level\cite{Cowan2011}. The fitted signal strength is
constrained to $\mu\geq0$, with a numerical search ceiling of 20. If the
lower endpoint reaches zero it is retained as a boundary endpoint; failure
to find an upper crossing is recorded rather than assigned a finite width.
This is a bounded profile interval with an asymptotic threshold, not a
pseudo-experiment-calibrated confidence belt.

The primary descriptive metric is the nominal Asimov width
\begin{equation}
 \W(S,k)=\mu_{\mathrm{up}}(S,k)-\mu_{\mathrm{low}}(S,k),
 \label{eq:width}
\end{equation}
at injected $\mu=1$ for subset $S$ and training seed $k$. Each candidate's
Asimov observation equals its own nominal expected count vector, including
the auxiliary expectations. Thus no Poisson fluctuations enter the nominal
comparison. An Asimov width measures the expected resolution under that
model; whether the corresponding interval has the advertised coverage is a
separate question, especially at small yields or near the physical boundary.

\subsection{Coalition values and paired summaries}

We define the coalition value by $v_k(S)=-\W(S,k)$, so larger values mean
narrower intervals. For group $G\in N$, its Shapley contribution is
\begin{equation}
 \phi_{G,k}=\sum_{S\subseteq N\setminus\{G\}}
 \frac{|S|!\,(3-|S|)!}{4!}
 \left[v_k(S\cup\{G\})-v_k(S)\right].
 \label{eq:shapley}
\end{equation}
The complete sixteen-coalition vector is available for every seed, so the
sum is evaluated exactly, without sampling feature orderings. It obeys
\begin{equation}
 \sum_{G\in N}\phi_{G,k}=v_k(N)-v_k(\varnothing).
 \label{eq:efficiency}
\end{equation}
The contributions have units of signal-strength interval width. They are
neither percentages of physical information nor coefficients of an additive
model for the event distribution.

Complementarity is examined through conditional second differences,
\begin{align}
 I_{GH\mid S,k}={}&v_k(S\cup\{G,H\})-v_k(S\cup\{G\})\nonumber\\
                 &-v_k(S\cup\{H\})+v_k(S),
 \label{eq:interaction}
\end{align}
where $S\subseteq N\setminus\{G,H\}$. There are six group pairs and four
conditioning subsets for each pair, giving 24 contrasts per seed. A positive
value means that their joint width reduction exceeds the sum of their
separate reductions conditional on $S$. This metric-specific interaction
does not imply causal synergy or a positive mutual-information contribution.

For two nonempty subsets we also retain direct paired reductions,
\begin{equation}
 R_k(S,T)=1-\frac{\W(S,k)}{\W(T,k)}.
 \label{eq:relative}
\end{equation}
All 105 unordered nonempty-subset pairs are available. Each contrast is
calculated within a seed and only then summarized across seeds by its
median. Taking a ratio of marginal medians would generally answer a different
question. Similarly, group contributions are computed before their medians;
the four median contributions need not satisfy Eq.~\eqref{eq:efficiency},
which is a per-seed identity.

The plots show all five seed values and their median. The broader run also
enumerates all $5^5=3125$ ordered resamples of the complete seed vectors to
describe stability of median summaries. Such resampling is conditional on
the shared MC and training prescription. It does not turn five initializations
into five independent datasets or supply a total confidence interval. BC and
AC are highlighted because of the observed exploratory ranking; they are not
retrospectively described as independently confirmed, prespecified winners.

\section{Results}
\label{sec:results}

\subsection{Compact inputs and expected precision}

Figure~\ref{fig:widths} displays the nominal widths for every nonempty
combination. The constant reference gives $\W=\WidthEmpty$ for each seed.
The lowest median width is obtained by BC, $\WidthBC$, closely followed by
AC, $\WidthAC$. Relative to the reference, the corresponding median
reductions are \GainBC\% and \GainAC\%. The full nineteen-variable ABCD
input has median width \WidthABCD{}, corresponding to a \GainABCD\%
reduction. All nominal coalition vectors are complete, so these comparisons
do not exclude unsuccessful seeds or substitute values for missing models.

\begin{figure*}[tbp]
\includegraphics[width=\textwidth]{subset_widths}
\caption{Nominal T1 Asimov interval widths at $\mu=1$. Combinations are ordered
by their median width. Coloured points show the five training seeds, and black
ticks show medians. The dashed line is the constant-score reference
$\Mempty$. Vertical offsets separate seed points only; they carry no
statistical meaning. The common mass grid has one bin from 105 to 140 GeV.
The seed spread is not a finite-MC or total confidence interval.}
\label{fig:widths}
\end{figure*}

The comparison with the full input is consistent in sign across the five
seeds. Both BC and AC yield narrower intervals than ABCD in every same-seed
comparison, with median reductions of \PairGainBC\% and \PairGainAC\%,
respectively [Fig.~\ref{fig:compact}]. Their compactness is substantial:
BC uses six inputs and AC ten, compared with nineteen for ABCD. Nevertheless,
the observed reduction in interval width relative to ABCD is modest. The
agreement of signs across seeds is useful descriptive evidence about this
training procedure, but cannot be assigned the evidential meaning of five
independent MC replications.

\begin{figure}[tbp]
\includegraphics[width=\columnwidth]{compact_comparisons}
\caption{Within-seed reductions in nominal interval width for BC and AC
relative to ABCD, computed using Eq.~\eqref{eq:relative}. Connecting lines
guide the eye between the discrete seeds. Both reductions are positive for
all five seeds. These comparisons were selected for discussion after the
complete subset ranking was examined.}
\label{fig:compact}
\end{figure}

BC's leading median should not be interpreted as universal first place.
AC gives the smallest width at seeds 42 and 46, BC at seeds 43 and 44, and
BCD at seed 45. The AC and BC distributions overlap, and AC wins their
direct comparison in two of five seeds. Thus the stable feature of the
result is the competitiveness of several compact combinations, rather than
a uniquely resolved ordering between the two leading choices. The complete
per-seed values in Appendix~\ref{app:nominal} make this distinction visible
without relying on rounded ranks.

The low yield scale also places the percentage changes in perspective.
Even the narrowest median interval has a width around 1.51 in a parameter
whose injected value is one. A reduction of roughly nine percent relative
to the constant count is a useful comparison of two procedures in the same
model; it is not a precision measurement of the Higgs signal strength. The
large interval scale and coarse mass template motivate retaining the
absolute widths alongside relative improvements.

\subsection{Contributions and conditional interactions}

The median Shapley contributions are
\begin{align}
 &(\phi_A,\phi_B,\phi_C,\phi_D)_{\mathrm{med}}\nonumber\\
 &\hspace{8mm}= (\PhiA,\PhiB,\PhiC,\PhiD).
 \label{eq:observedphi}
\end{align}
Groups B and C contribute most strongly and have comparable median values.
Group A contributes positively but less, whereas D is negative in every
seed [Fig.~\ref{fig:shapley}]. Exact recomputation from the published
coalition widths reproduces the stored contributions and their per-seed
efficiency identity. This algebraic consistency verifies the attribution
calculation, without adding information about its finite-MC uncertainty.

\begin{figure}[tbp]
\includegraphics[width=\columnwidth]{shapley}
\caption{Exact Shapley contributions to negative interval width. Colours
identify the same seeds as in Fig.~\ref{fig:widths}; black ticks denote
medians. Positive values indicate an average reduction of width when a
group is added over all possible conditioning coalitions. D is negative in
each seed. The points describe variability across fitted networks on the
same MC, rather than total uncertainty or intrinsic physical information.}
\label{fig:shapley}
\end{figure}

The second differences refine this picture. Without a conditioning group,
the AC interaction is positive in every seed, with median approximately
$+0.04246$ in width units. The corresponding AB interaction is negative
in every seed, with median $-0.05344$. Under this value function, A and C
provide a joint gain larger than the sum of their separate gains, while
A and B exhibit diminishing combined gain. Other conditioning coalitions
can change the numerical size or sign of a contrast. All 24 contrasts,
rather than only these two examples, are listed in
Appendix~\ref{app:interactions}.

This pattern has a plausible representation-level interpretation. Dilepton
masses and separations expose structures that a small learner would
otherwise have to infer indirectly, while collective kinematics can provide
different discrimination or acceptance information. Overlap among the
representations can make the gain from adding one depend strongly on what
is already present. However, the numerical attribution cannot determine
which physical mechanism dominates. Establishing that mechanism would
require matched controls that separate production information, mass
dependence, reconstruction effects, and learning convenience.

In particular, the negative D contribution is not evidence against the
physical relevance of decay angles. It states that adding this angular
representation tends to widen the interval after retraining and rebuilding
the prescribed score categories and finite-statistics likelihood. The
comparison combines effects of estimation, optimization, category support,
and the input-dependent parameter count. An ideal information-preserving
analysis could disregard unhelpful coordinates, whereas this finite
training procedure need not do so. A negative contribution is therefore
compatible with physically informative observables.

\subsection{Classification and inference rankings}

Figure~\ref{fig:auc} compares the validation AUC with the expected interval
width. There is a broad association: more discriminating classifiers
usually yield narrower intervals in this family. Across the fifteen
candidate medians, the descriptive Pearson and Spearman correlations are
approximately $-0.945$ and $-0.911$, respectively. These correlated,
exploratorily examined candidates are not independent observations for a
population-level significance test, and no such test is assigned to the
correlation coefficients.

\begin{figure}[tbp]
\includegraphics[width=\columnwidth]{auc_width}
\caption{Validation absolute-weight AUC versus nominal T1 interval width.
Faint coloured points are individual seed results; black points are the
separate AUC and width medians for each nonempty combination. Selected
combinations are labelled for readability. The overall association is
strong, but the rankings differ: AC has the highest median AUC and BC the
smallest median width. The constant baseline is omitted because its AUC is
not recorded as a trained-model validation result.}
\label{fig:auc}
\end{figure}

The strongest classification ranking nevertheless differs from the
strongest inference ranking. AC has the highest median AUC, while BC has
the smallest median width. ABCD has a higher median AUC than BC but a wider
interval. This is expected to be possible because AUC does not include
the signal/background normalization, the two-category compression, or the
penalty from finite template statistics. A change that improves ordering
over much of the score range can leave the particular threshold used by
the likelihood less favourable.

The distinction is also sharpened by the weighting convention. Validation
AUC is evaluated with absolute weights, whereas the category expectations
are signed physical yields. Thus even before compression the evaluation
measures are not identical. The study supports evaluating representations
through the intended inference pipeline, but does not compare conventional
training with an inference-optimized learner. It supplies no benchmark
against INFERNO, a matrix-element likelihood, or an optimal statistic.

\section{Validation and limitations}
\label{sec:validation}

\subsection{Finite-MC resampling}

The event-MC bootstrap resamples physical groups in the calibration and
template roles separately, using common draws across all candidates. The
networks remain fixed; score thresholds are refitted, category templates
are rebuilt, and the complete inference and attribution calculation is
repeated on the frozen nominal mass grid. The procedure therefore addresses
the finite calibration/template sample conditional on the trained models.
It does not include retraining, assessment-parent uncertainty, or variation
of the underlying generator and detector model.

The prescribed ensemble contains 200 replicas. Of these, 161 are complete
and valid and 39 fail the statistical-support requirements. The registered
rule requires all 200 complete replicas for formal percentile intervals,
so no such intervals are reported. Summaries conditional on the successful
replicas would describe a selected subset of the resampling distribution.
They cannot be silently substituted for the missing unconditional
bootstrap result. The failure fraction, 19.5\%, is itself evidence that
the present finite-MC support is a material limitation.

Consequently, the nominal differences and the seed patterns do not yet
establish a gain that survives finite-simulation uncertainty. The missing
intervals cannot be repaired by combining seed spread with a conditional
bootstrap spread or by generating additional random replicas until all
desired results appear. More draws from the same parent improve Monte
Carlo integration of a resampling procedure; they do not supply absent
support or additional independent simulated events.

\subsection{Pseudo-experiments and repeated calibration}

Model-self and assessment-parent pseudo-experiments are generated at
$\mu=0,1,2$, with 500 trials per candidate and injection for each of the five
training-seed blocks. The first tests closure under the fitted model; the
second uses the frozen assessment parent to examine a distinct source of
expectations. Poisson auxiliary observations are regenerated according to
the registered T1 prescription. Each block contains sixteen candidates. The retained report
contains valid numerical outputs for all fifteen model-self cells and all
fifteen assessment cells. Numerical completion is distinguished from the
physical applicability and independence of those diagnostics.

Within a block, candidate count vectors are coupled by sharing total
Poisson draws and using common uniforms for monotone category allocation.
This construction preserves each candidate's marginal Poisson law. It does
not reproduce the covariance of classifications of a common set of physical
events. As a result, paired errors describe the registered artificial
common-random-number coupling, not physical cross-model covariance.
Matching trial indices in different training-seed blocks has no event-pairing
meaning. A sensitivity calculation with independent category allocation is
pending, so coupled pairwise diagnostic uncertainties are not used as a
primary reliability claim.

Repeated-calibration diagnostics, denoted T2, use twenty outer
calibration-group bootstrap replicas per seed and one hundred inner
pseudo-experiments at $\mu=1$ for each replica. Each replica's thresholds
are applied consistently to the template and pseudo-observation
construction, with the nominal mass grid held fixed. Networks and the
template/assessment parent events are not regenerated. All five T2 cells
have valid numerical outputs. The outer calibration replicas are the
procedure-level units; the inner trials are conditional simulations within
each unit. Treating the resulting 2,000 fits as independent replications
of the entire analysis would overstate the evidence.

Table~\ref{tab:validation} gives conditional coverage for the three
highlighted candidates. Entries are medians over training seeds, not
pooled binomial estimates. At the nominal 68\% level, assessment coverage
ranges from 0.634 for AC and ABCD to 0.668 for BC. The corresponding T2
values are 0.648, 0.650, and 0.666. At 95\%, the displayed fractions are
near or somewhat above nominal. These values indicate that the interval
construction deserves further coverage study; they do not support an
unqualified statement that the narrowest interval is the most reliable.

\begin{table*}[tbp]
\caption{Conditional coverage at $\mu=1$ and the status of the validation
layers. Numerical entries are five-seed medians of the reported conditional
coverage. T2 averages inner conditional coverage within each outer
calibration procedure before the seed summary. These are not pooled
confidence intervals or independent validation results.}
\label{tab:validation}
\begin{ruledtabular}
\begin{tabular}{lccp{8.0cm}}
Diagnostic / candidate & 68\% & 95\% & Interpretation or remaining limitation \\
\hline
Assessment: BC & \AssessmentBCSixtyEight & \AssessmentBCNinetyFive & Frozen assessment parent; non-independent access review. \\
Assessment: AC & \AssessmentACSixtyEight & \AssessmentACNinetyFive & Same parent convention and artificial within-seed coupling. \\
Assessment: ABCD & \AssessmentABCDSixtyEight & \AssessmentABCDNinetyFive & No total or externally validated coverage claim. \\
T2: BC & \TtwoBCSixtyEight & \TtwoBCNinetyFive & Twenty outer calibration replicas; fixed networks and parents. \\
T2: AC & \TtwoACSixtyEight & \TtwoACNinetyFive & One hundred conditional inner trials per outer replica. \\
T2: ABCD & \TtwoABCDSixtyEight & \TtwoABCDNinetyFive & Outer procedures must not be replaced by pooled inner trials. \\
\hline
Nominal comparison & \multicolumn{2}{c}{80/80 valid} & Fixed-model Asimov comparison; five complete coalition vectors. \\
Model-self / assessment & \multicolumn{2}{c}{15/15 cells each} & Conditional diagnostics; physical-event covariance not reproduced. \\
Event-MC bootstrap & \multicolumn{2}{c}{161/200 valid} & Formal percentile intervals unavailable; 39 support failures. \\
Independent validation & \multicolumn{2}{c}{Pending} & Signed-MC approximation, access history, and physical variations. \\
\end{tabular}
\end{ruledtabular}
\end{table*}

For a single set of 500 independent conditional Bernoulli outcomes, the
binomial standard error at coverage 0.68 would be about 0.021. This provides
a scale for simulation noise within a cell, but cannot be applied directly
to a median across shared-MC seeds or to nested T2 trials. The report
distinguishes successful-fit conditional coverage, the planned-denominator
fraction of trials that both succeed and cover, and the fit-failure rate.
These distinctions matter whenever fits fail. At $\mu=0$, the physical
boundary additionally alters the behaviour of the intervals; conservative
boundary coverage is not evidence of validated coverage away from zero.

\subsection{Scope of the physical interpretation}

Three limitations constrain the interpretation beyond numerical precision.
First, the comparison is specific to a restricted simulated process model
and a single reconstructed final state. Generator composition, normalization,
negative-weight origin, radiation conventions, and the detector description
require independent source-level validation. The use of an open simulated
collection does not by itself provide that validation. Nor does a finite-MC
nuisance model represent the full set of theoretical and experimental
systematic uncertainties in a Higgs measurement.

Second, the learned representations are compared at one training budget
with one architecture family. Their input dimensions and first-layer
parameter counts differ. The results therefore cannot establish that BC
or AC is sufficient for the underlying physics, or that it will dominate
with more simulation, a different learner, or a more informative likelihood.
A matrix-element reference, controlled capacity comparisons, and training
sample-size studies would help separate representation convenience from
missing information. Those experiments have not been performed as part of
the result reported here.

Third, the compact candidates were highlighted after observing the complete
combination ranking. Model-selection criteria themselves fluctuate, and
choosing a favourable candidate can bias a subsequent interpretation of its
performance\cite{Cawley2010}. The five seeds and exhaustive contrasts improve
transparency, but do not eliminate this selection effect or create a
simultaneous-coverage guarantee for the 105 pair comparisons. Independent
confirmation requires a separately fixed candidate comparison and eligible
data or simulation whose previous use does not compromise that confirmation.
Relabelling the same events or merely changing a seed would not meet that
requirement.

These limitations leave a useful, narrower conclusion. In the executed
single-mass-bin analysis, conventional classifier training can yield compact
kinematic representations whose score partitions have smaller nominal
signal-strength intervals than the full input. Exact coalition accounting
helps identify where this occurs and shows why the contribution of one
group depends on the others. Whether the observed improvement persists
under adequate finite-MC support, validated signed-template modelling, a
resolved mass fit, and independently sourced variations remains open.

\section{Conclusions}
\label{sec:conclusions}

We have examined kinematic feature groups through their effect on the
expected precision of signal-strength inference in a controlled simulated
$H\to ZZ^*\to2e2\mu$ population. All fifteen nonempty combinations of
nineteen observables were evaluated with five trained networks each and a
deterministic empty-coalition reference. Explicit four-lepton mass was
excluded from the classifiers. The final T1 likelihood used a common,
frozen grid that had merged to one mass bin, so the comparison concerns
score-category information within the selected mass window.

The compact BC and AC inputs have median nominal widths of \WidthBC{} and
\WidthAC{}, compared with \WidthABCD{} for ABCD. They improve over the
constant reference by \GainBC\% and \GainAC\% and outperform ABCD in all
five same-seed comparisons. Their relative ordering is not stable enough
to identify a universal winner. Exact Shapley attribution gives the largest
positive contributions to B and C, a smaller positive contribution to A,
and a negative contribution to D for this particular procedure. The AUC
and interval-width rankings differ despite their strong overall association.

The numerical result remains exploratory. Event-group bootstrap failures
prevent formal finite-MC percentile intervals, pseudo-experiment pairing is
artificial, and independent physical and assessment-history validation is
incomplete. The negative angular contribution does not imply absent angular
information, and the nominal reductions do not establish a precision gain
over a resolved mass-spectrum analysis. The complete coalition comparison
instead defines concrete compact candidates and validation questions for a
subsequent independently controlled study.

\begin{acknowledgments}
[Funding and acknowledgments to be supplied by the authors.]
\end{acknowledgments}

\section*{Data and computational materials}
Only simulated members of the cited public collection were used. The
accompanying draft materials contain the aggregate result snapshot, its
source-file checksums, plotting code, and compilation instructions. The
snapshot records the existing frozen analysis; preparing this manuscript
did not retrain models, refit likelihoods, or generate new pseudo-experiments.
The run provenance and the source revision used for manuscript verification
are distinguished in the accompanying evidence record. A permanent archival
software and result release has not been established by this draft.

\clearpage
\onecolumngrid
\appendix
\section{Complete nominal comparison}
\label{app:nominal}
Table~\ref{tab:nominal} supplies the full nominal coalition vector for each
seed. Displayed values are rounded to five decimal places; ranking and
contrasts use the full-precision aggregate snapshot. The AUC column is the
median of the selected-checkpoint validation values under absolute physical
weights. The blank AUC for the deterministic reference is preserved rather
than replaced with an assumed trained-model measurement.

\begin{table}[!h]
\caption{Nominal $\W$ by training seed, its median, and median validation AUC.
The empty coalition is $\varnothing$. The five seed columns share the same
underlying MC. Their spread is not a total uncertainty estimate.}
\label{tab:nominal}
\small
\centering
\begin{ruledtabular}
\begin{tabular}{lrrrrrrr}
Subset & 42 & 43 & 44 & 45 & 46 & Median $\W$ & Median AUC \\
\hline
\input{generated/nominal_rows}
\end{tabular}
\end{ruledtabular}
\end{table}

\clearpage
\section{Complete conditional interactions}
\label{app:interactions}
Table~\ref{tab:interactions} reports every second difference in
Eq.~\eqref{eq:interaction}. The table specifies the conditioning coalition
because a group-pair label alone does not define an interaction. Positive
and negative values refer to the negative-width coalition value and cannot
be converted into percentages of information. These descriptive contrasts
do not have a simultaneous-confidence interpretation.

\begin{table}[!h]
\caption{All 24 conditional interactions in $\W$ units. Each row identifies
the added pair $GH$ and conditioning subset $S$. The last column is the median
of the five within-seed contrasts, rather than a contrast of median widths.}
\label{tab:interactions}
\small
\centering
\begin{ruledtabular}
\begin{tabular}{llrrrrrr}
$GH$ & $S$ & 42 & 43 & 44 & 45 & 46 & Median \\
\hline
\input{generated/interaction_rows}
\end{tabular}
\end{ruledtabular}
\end{table}

\clearpage
\twocolumngrid
\bibliography{references}
\end{document}
